The experiments show a clear gap between the \emph{theoretical} speed-up of Grover’s algorithm and its \emph{practical} usefulness on student-accessible quantum hardware. In noise-free simulations, Grover amplification behaves as expected for all three problems, confirming that the encodings, oracles, and diffusion operators are logically correct. On real hardware, however, performance is dominated by circuit depth and two-qubit gate errors, so the theoretical advantage is usually lost before a useful number of Grover iterations can be executed.
\\
\\
A main reason Grover underperforms for small instances is constant overhead. Turning a classical constraint check into a reversible oracle requires extra ancilla qubits, multi-controlled operations, and uncomputation. For the instance sizes that fit on hardware, this overhead can be comparable to (or exceed) the full runtime of classical solvers. In addition, all three problems have an exploitable structure. Sudoku can be solved efficiently by classical backtracking with pruning, and Lights Out can be solved via linear algebra. WFC uses local constraint propagation and collapse, which produces valid maps quickly in practice. Grover’s quadratic improvement applies to \emph{unstructured brute-force search}, but the strongest classical baselines are \emph{structure-aware}, which is why classical methods dominate for the tested instance sizes.
\\
\\
Across all three games, the most important practical limiter is circuit depth after transpilation. Limited device connectivity introduces SWAP gates and increases two-qubit gate volume, which rapidly accumulates error. This effect is visible in Lights Out, where the circuit grows from 156 logical gates to 3287 after transpilation, and the measured solution frequency drops from $\approx 96.13\%$ (noiseless simulator) to $\approx 37.61\%$ (noise model) and $\approx 16.42\%$ (hardware). In Sudoku, the $2\times 2$ instance still demonstrates Grover successfully on hardware, but the partially filled $4\times 4$ case fails despite sufficient qubits because the oracle is deep and requires many Grover iterations; compilation overhead and noise flatten the distribution. In WFC, the $3\times 3$ simulator run achieves $99.6\%$ valid outputs (1020/1024), yet the IBM hardware run yields only $\approx 1.4\%$ valid outputs (14/1024), consistent with the transpiled depth rising to $\mathcal{O}(10^5)$ and overwhelming any amplification.
\\
\\
Ancilla qubits create a trade-off that affects all three implementations. More ancillae can reduce depth by enabling more efficient decompositions of multi-controlled gates, but they also increase routing complexity and can increase SWAP overhead. Therefore, reducing qubits can sometimes help (less routing) and sometimes hurt (deeper decompositions and more sequential operations). In practice, the best choice depends on the target backend and circuit layout, not only on the algorithmic design.
\\
\\
Looking forward, improvements in (i) two-qubit error rates, (ii) coherence times, and (iii) connectivity would directly increase the maximum usable circuit depth and allow more Grover iterations before the signal is lost. Error correction would change the picture more fundamentally, but it requires large overhead beyond current student-accessible devices. On the algorithm side, progress would come from shallower oracles (better constraint encodings, improved MCX synthesis, and ancilla reuse) and hybrid approaches where classical pruning reduces the search space before applying Grover. Even with better hardware, highly structured problems are likely to remain classical-favorable unless the quantum approach also exploits structure rather than only accelerating brute-force search.

